\documentclass[12pt]{article}

% You need graphics package
\usepackage{color}
\usepackage{graphicx}
\usepackage{tabularx}

\usepackage{pst-all}

% \usepackage[annot]{advi}
\usepackage{hyperref}

%% \usepackage[annot]{advi} 
%% must occur after packages using PStricks macros 
%% (it redefines some of them) 

\definecolor{red}{rgb}{1.0,0.0,0.0}
\definecolor{blue}{rgb}{0.0,0.0,1.0}

\def\email#1{\textcolor{gray}{$<${#1}$>$ }}

\def\key#1{\textcolor{red}{#1}}
\def\ikey#1#2#3{\key{#1} & #2 & -- & #3\\}
\def\arg{\textcolor{blue}{arg }}

\pagestyle{empty}

\def\advifooter{\vbox to 0em{\vbox to \vsize {\vfill
\center{Press: \key{n}ext page \key{p}revious page \key{q}uit \vss}}}}

\def\adviheader{\noindent
{\bf\Large Active dvi}\\
\includegraphics[width=\textwidth]{bar.jpg.eps}}

\let \Newpage \newpage
\def \newpage {\Newpage \advifooter\adviheader}

\def\adviemptyfooter{\vbox to 0em{\vbox to \vsize {\vfill
\hfill Press \key{q} to quit \hfill~} \vss}}

\def\lastpage{\Newpage\adviemptyfooter\adviheader}

\begin{document}
 
\newpage

%\section*{Active dvi}

\medskip

\begin{tabularx}{\linewidth}{Xl}
The presentation tool of choice for discriminating hackers.
See \texttt{http://pauillac.inria.fr/advi/}
&
\raisebox{-30pt}{\includegraphics[width=2cm]{advilogo.eps}}
\end{tabularx}

\subsection*{Usage}
\texttt{advi [OPTIONS] DVIFILE}\\
Try \texttt{advi -help} for more information.

\subsection*{Key bindings}

Advi recognizes the keystrokes listed below when typed in its window.
Some keystrokes may optionally be preceded by a number, called \arg below, whose
interpretation is keystroke dependant. If \arg is unset, its value is
1, unless specified otherwise.

Advi maintains an history of previously visited pages organized as a stack. 
Additionnally, the history contains marked pages which are stronger than
unmarked pages. 

\vspace{\stretch{1}}

\noindent
\begin{tabularx}{\linewidth}{clcX}
\ikey{?}{info}{Quick info and key bindings help.}
\ikey{q}{quit}{End of show.}
\ikey{space}{continue}
{Move \arg pauses forward if any, or do as \key{return} otherwise.}
\ikey{n}{next}
{Move \arg physical pages forward, leaving the history unchanged.}
\ikey{p}{previous}
{Move \arg physical pages backward, leaving the history unchanged.}
\ikey{g}{go}
{If \arg is unset move to the last page.
 If \arg is the current page do nothing.
 Otherwise, push the current page on the history as marked, and move
 to the physical pages \arg.}
\ikey{,}{begin}{Move to the first page.}
\ikey{.}{end}{Move to the last page.}
\end{tabularx}

\vspace{\stretch{1}}

\newpage

\vspace{\stretch{1}}

\noindent
\begin{tabularx}{\linewidth}{clcX}
\ikey{N}{next pause}{Move \arg pauses forward (equivalent to continue).}
\ikey{P}{previous pause}{Move \arg pauses backward.}
\ikey{$<$}{smaller}{Scale down the resolution by scalestep (default $\sqrt{\sqrt 2}$).}
\ikey{$>$}{bigger}{Scale up the resolution by scalestep (default $\sqrt{\sqrt 2}$).}
\ikey{c}{center}{Center the page in the window, and resets the default resolution.}
\ikey{\char35}{fullpage}{Remove margins around the page and change
the resolution accordingly.}
\ikey{\verb"^"L}{redisplay}{Redisplay the current page to the first
pause of the page.}
\ikey{r}{redraw}{Redraw the current page to the current pause.}
\ikey{R}{reload}{Reload the file and redraw the current page.}
\ikey{h}{page left}
{Moves one screen width toward the left of the page. Does nothing if the
  left part of the page is already displayed}
\ikey{l}{page right}
{Moves one screen width toward the right of the page. Does nothing if the
  right part of the page is already displayed}
\ikey{j}{page down}
{Moves one screen height toward the bottom of the page. Jumps to the top of
  next page, if there is one, and if the bottom of the page is already
  displayed.}
\ikey{k}{page up}
{Moves one screen height toward the top of the page. Jumps to the bottom
 previous page, if there is one, and if the top of the page is already
 displayed.}
\end{tabularx}

\vspace{\stretch{1}}

\newpage

\vspace{\stretch{1}}

\noindent
\begin{tabularx}{\linewidth}{clcX}
\ikey{return}{forward}
{Push the current page on the history stack, and move forward n physical pages.}
\ikey{tab}{mark and next}
{Push the current page on the history as marked, and move forward n physical pages.}
\ikey{backspace}{back}
{Move \arg pages backward according to the history. The history stack is poped, accordingly.}
\ikey{escape}{find mark}
{Move \arg marked pages backward according to the history.
 Do nothing if the history does no contain any marked page.}
\ikey{f}{load fonts}{Load all the fonts used in the documents.  By default, fonts are loaded only when needed.}
\ikey{F}{make fonts}
{Does the same as \key{f}, and precomputes the glyphs of all characters used in the document.
This takes more time than loading the fonts, but the pages are drawn faster.}
\ikey{C}{clear}{Erase the image cache.}
\end{tabularx}

\vspace{\stretch{1}}

\lastpage

\subsection*{Acknowledgments}

Authors and contributors are :

\begin{itemize}
\item Alexandre Miquel \email{Alexandre.Miquel@inria.fr}
\item Jun Furuse \email{Jun.Furuse@inria.fr}
\item Didier R\'emy \email{Didier.Remy@inria.fr}
\item Xavier Leroy \email{Xavier.Leroy@inria.fr}
\item Pierre Weis \email{Pierre.Weis@inria.fr}
\item Didier Le Botlan \email{Didier.Le\_Botlan@inria.fr}
\item Alan Schmitt \email{Alan.Schmitt@inria.fr}
\item Roberto Dicosmo \email{Roberto.Dicosmo@inria.fr}
\end{itemize}

\vspace{\stretch{1}}

\begin{center}
This program is distributed under the GNU LGPL. \\
See the enclosed file COPYING.
\end{center}

\vspace{\stretch{3}}

\end{document}
